\chapter{関連研究}
\label{chap:related}

野生動物の監視では，画像の確認作業を減らすこと，出現情報を速やかに伝えること，電源と通信の制約下で運用することが課題となる．
本章では，トレイルカメラ画像の自動解析，通信型の監視・通知，省電力化と通信量削減に関する研究を整理し，本研究の位置付けを述べる．

\section{トレイルカメラ画像の自動解析}
Norouzzadehら\cite{Norouzzadeh2018}は，約320万枚の画像を含むSnapshot Serengetiデータセットを対象に，深層学習による48種の動物の識別，個体数および行動の推定を検討している．
さらに，予測の信頼度に応じて自動処理する画像と人手で確認する画像を分けることで，大量の画像に対する確認作業の削減を示している．
この研究は，画像解析を人手による確認と組み合わせる意義を示すものである．
ただし，動物種の分類正解率は，物体検出の再現率や，通信前に空撮影画像を除外できる割合とは異なる指標である．

Tabakら\cite{Tabak2019}は，米国の5州で収集した約337万枚の画像を用いて，ResNet-18に基づく動物種分類モデルを学習した．
米国内の学習に用いなかった画像では98\%の分類正解率を報告している一方，カナダの別データに対しては82\%となっている．
この結果は，大量画像の自動分類の有効性を示すとともに，学習時と異なる地域や対象に対して性能が変化し得ることを示している．
したがって，他のデータセットでの高い性能を，設置先の画像にそのまま適用できるとは限らない．

Beeryら\cite{Beery2019}は，汎用的な動物検出器で空撮影画像を選別し，検出領域を切り出して対象地域に応じた分類器へ入力する処理を提案している．
動物の位置検出と種分類を分離することにより，画像全体を直接分類する場合に比べ，背景に依存する問題を軽減しながら確認作業を支援する考え方である．
この処理は，動物が含まれる可能性のある画像を先に抽出し，後段の処理対象を絞るという点で，本研究の段階的な画像選別に関係する．

Fennellら\cite{Fennell2022}は，カナダのカメラトラップ画像に対するMegaDetectorの検出結果を，人手による分類と比較している．
信頼度閾値を0.9とした条件で，動物画像に対する適合率82\%，再現率92\%を報告し，検出器を組み込むことによる目視確認作業の短縮を示している．
一方で，適合率と再現率がともに100\%ではないことから，汎用検出器の出力にも誤りが含まれる．
本研究においてMegaDetectorを再判定や参照評価に用いる場合にも，その出力を無条件に正解とみなさず，人手による確認との違いを明確にする必要がある．

これらの研究は，蓄積画像を効率よく解析するための基盤を与える．
本研究ではその知見を踏まえ，画像を遠隔側に送る前の選別と，送信後の再判定に役割を分け，画像確認だけでなく通信と通知を含めたシステムとして検討する．

\section{通信型の監視・通知システム}
Whytockら\cite{Whytock2023}は，カメラトラップにAIによる画像判定とIridium衛星通信を組み合わせ，判定ラベルとメタデータを利用者に送るシステムを開発している．
オランダでの試験とガボンの森林での試験を行い，撮影から通知受信までの時間の中央値として7.35分を報告している．
一方で，設置状態が適切でなく衛星ネットワークに接続できない場合には，長い通知遅延も生じている．
この研究は，携帯電話網への依存を避ける通信手段と即時的な情報提供の可能性を示しているが，画像そのものを送る本研究の構成とは，通知内容と通信条件が異なる．

Patkarら\cite{Patkar2025}は，PIRセンサ，Raspberry Piおよび独自に学習したYOLOv8nを組み合わせた野生動物侵入検知システムを提案している．
同システムでは，動物を検出した結果をTwilioによるSMS通知につなげており，撮影地点側での推論を獣害対策の通知に利用している．
このような構成では，通知までの処理を現地の計算機で実行できる一方，監視地点ごとの計算資源と電源をどのように確保するかが設計上の検討事項となる．

Zualkernanらは，エッジでの動物種分類とクラウド上の情報管理を組み合わせるIoT構成を提案している\cite{Zualkernan2020}．
2022年の論文\cite{Zualkernan2022}では，エッジ端末で画像を分類し，ラベル，時刻および端末識別情報をLoRaWAN経由で送信して，Firebaseとモバイルアプリケーションで管理するシステムを示している．
また，Raspberry Pi，Google CoralおよびJetson Nanoでの実行を比較し，推論時間と消費電流に加えて，端末向けのモデル変換に伴う分類性能の変化を検討している．
この研究は，画像認識と通信を組み合わせて遠隔地へ出現情報を伝える先行例である．
一方，広域側へ送るのは主に分類結果であり，本研究が一次選別後の画像を転送し，クラウドで再判定する構成とは異なる．
画像転送には多くのデータ量が必要となるが，遠隔側で画像を再解析し，利用者が判定の根拠を確認できるという違いがある．

通信型の監視では，検出器の精度だけでなく，検知から通知までの経路全体を評価する必要がある．
また，ラベルのみを送る構成，画像を送る構成，利用者が画像を閲覧する構成では目的が異なるため，通知時間や通信量を比較する際には，送信内容と測定区間をそろえる必要がある．

\section{省電力化と通信量削減}
Cunhaら\cite{Cunha2021}は，組込み機器上で空撮影画像を除外するため，画像分類器と物体検出器の選別性能および推論時間を比較している．
入力解像度の変更や量子化などを検討し，動物を含む画像の再現率を96\%とするように信頼度閾値を調整した条件で，空画像を除外する性能を評価している．
この評価方法は，動物画像の見逃しを一定の範囲に抑えながら，転送対象をどこまで減らせるかを検討する本研究に関係する．
ただし，空画像のうち正しく除外した割合と，全撮影画像に対する転送枚数の削減率は分母が異なるため，両者を区別する必要がある．

また，Zualkernanら\cite{Zualkernan2022}の機器比較は，推論の高速化だけで端末の選択を決められないことを示している．
機器ごとの消費電流を比較する場合には供給電圧と処理時間を併せて確認し，待機を含む運用時の電力量と分けて検討する必要がある．
本研究では，推論機器を各撮影地点に置く代わりに共有エッジサーバへ集約し，カメラ側の間欠動作と，集約側の常時稼働を合わせて評価する．

Saitoら\cite{IEICE2020}は，深層学習による野生動物検出と通知を行う，電池駆動の検出ノードを提案している．
単一の小型コンピュータボードを用いる構成と，二つのボードを用いる構成を扱い，対象動物が検出されない場合の処理を省くことなどによって省電力化を図っている．
検出だけでなく電源制約を考慮して装置を設計し，消費電力や通信成功率を実環境で評価している点が，本研究に関係する．

Satoら\cite{Sato2022}は，動体センサの感度調整，センサの検知範囲を制限する覆いの装着，画像のフレーム差分，および親機と子機への機能分離を提案している．
これらは，不要な検知，撮影，推論や通信を減らし，動作時と待機時の電力量を抑えるための手法である．
評価対象の装置では，複数の手法を組み合わせることで，基準構成に対する電力量の削減を示している．
ここから，通信回線の集約だけでなく，どの処理を各子機に残すかを検討する必要があることが分かる．

Monburinonら\cite{Monburinon2019}は，Raspberry Piなどのゲートウェイ機器で主要な画像認識を行う階層型エッジコンピューティングを提案している．
設置環境に応じた対象クラスを学習し，ゲートウェイ側で判定する構成を扱っている．
したがって，現地側に画像処理を配置する考え方や，ゲートウェイを利用する階層構成そのものは既に研究されている．
本研究では，そのような処理配置を踏まえ，撮影専用カメラ，共有エッジサーバ，クラウドでの再判定の役割を明示して設計する．

Velasco-Monteroら\cite{Velasco2024}は，カメラトラップへのAI処理の組込みを，ハードウェア，ソフトウェアおよび学習手法を合わせて検討している．
試作機と現地で収集したデータを用い，継続学習などによって設置環境に応じた判定性能の改善を図っている．
この研究は，端末内での画像選別と現地環境への適応を重視するものであり，本研究が推論を共有エッジサーバ側へ分離する設計とは，計算資源の配置が異なる．
どちらの配置が適切かは，カメラ台数や画像発生頻度，電源，通信条件と合わせて判断する必要がある．

Rahmanら\cite{Computers2025}は，YOLOv8mによる検出に，Kalmanフィルタと画像の外観特徴に基づく対応付けを組み合わせ，動物を追跡する処理を提案している．
同じ動物が連続フレームに現れる場合の識別を維持し，重複した情報の転送を減らすことを目的としている．
このような重複情報の削減は，動物が写っていない画像を除く空撮影画像の選別とは異なるが，いずれもエッジ側で送信対象を絞る手法である．
本研究で撮影サイクルを単位として画像を管理する際にも，空撮影の除外と，同一の出現に関する重複情報の削減を区別することが重要となる．

\section{本研究の位置付け}
以上の研究から，画像認識による確認作業の支援，現地側での推論と遠隔通知，機能分離による省電力化は，それぞれ既に検討されていることが分かる．
表\ref{tab:related-position}に，本研究と関連する設計上の観点を整理する．
各研究は対象データや通知内容が異なるため，表は同一条件での性能比較ではなく，処理の配置と目的の比較である．

\begin{table}[tbp]
  \centering
  \small
  \caption{関連研究と本研究の着眼点}
  \label{tab:related-position}
  \begin{tabularx}{\linewidth}{p{0.23\linewidth}XX}
    \toprule
    研究 & 主な構成・処理 & 本研究との関係 \\
    \midrule
    Beeryら\cite{Beery2019} & 動物検出と種分類を分離した画像解析 & 段階的な画像選別の基礎 \\
    Whytockら\cite{Whytock2023} & 現地判定と衛星回線によるラベル通知 & 遠隔通知に着目．転送内容と通信経路は異なる \\
    Zualkernanら\cite{Zualkernan2022} & エッジ分類とLoRaWANによるラベル送信 & 画像転送・クラウド再判定とは送信内容が異なる \\
    Cunhaら\cite{Cunha2021} & 組込み機器での空画像除外と推論時間の比較 & 見逃しを抑えながら送信対象を減らす評価の基礎 \\
    Satoら\cite{Sato2022} & 不要な処理の抑制と親機・子機への機能分離 & 各機器の役割と電力量を合わせて検討 \\
    Rahmanら\cite{Computers2025} & 動物追跡による重複転送の削減 & 空撮影の選別とは異なる送信量削減 \\
    本研究 & 撮影専用カメラ，共有エッジでの一次選別，クラウドでの再判定 & 複数カメラの画像集約と二段階判定を組み合わせて評価 \\
    \bottomrule
  \end{tabularx}
\end{table}

本研究の設計上の特徴は，複数配置するカメラを画像取得に特化させ，画像の集約と一次選別を共有エッジサーバに担わせることである．
さらに，一次選別を通過した画像をクラウドへ送り，別の検出モデルで再判定して通知の対象を決める．
これにより，カメラ個別の広域通信回線を必要としない構成と，広域側へ送る画像を選別する構成を組み合わせる．

ただし，段階的な処理の導入だけで，検出性能，消費電力，通知時間が同時に改善されるとは限らない．
一次選別で除外した動物画像は後段の再判定で回復できず，エッジサーバの常時稼働には電力が必要となる．
このため，本研究では，画像の削減と動物の見逃し，通信の到達性と通知時間，各ノードの消費電力を対応付けて評価する．
低コスト化や長期間稼働は設計の目標として扱い，達成の範囲は具体的な実装条件と評価結果に基づいて議論する．
